Fix a calibrated \(\kappa\) and \(\gamma>d\). By \(\mathrm{def:transverse-model}\), \(\mathcal M_{\mathrm{tr}}(\gamma;\kappa)=\mathcal M(1,\gamma;\kappa)\). The class is nonempty by \(\mathrm{thm:calibrated-witness}\), and \(\mathrm{thm:identified-version}\) shows that its causal upper set is an observed-law functional modulo null sets and that the loss is well defined.

For the upper bound, \(\mathrm{lem:weighted-margin-transfer}\) gives, for all sufficiently large \(n\),
\[
\sup_{P\in\mathcal M_{\mathrm{tr}}(\gamma;\kappa)}
 \mathbb E_P\!\left[d_H(\widehat G_n^{\mathrm{LP}},G_P;P)\right]
 \le C_1(r_n^*)^2.
\]
The proof of that lemma uses the exponential pointwise inequality for the exact boundary-adapted, spectrally clipped statistic in \(\mathrm{def:boundary-adapted-estimator}\); it does not infer a tail bound from a root-mean-square rate. Taking the infimum over all measurable set estimators yields
\[
\mathfrak R_n(1,\gamma;\kappa)\le C_1(r_n^*)^2,
\]
and the displayed estimator attains this inequality in both tuning branches.

For the converse, \(\mathrm{lem:calibrated-transverse-assouad-converse}\) constructs laws inside this same class. Its direct adjacent experiment gives the high-smoothness branch. In the low-smoothness branch its fuzzy adjacent experiment is controlled by \(\mathrm{lem:conditional-component-fuzzy-hellinger}\), whose calculation conditions on the common uniform design and factors the normalized conditional laws over support components. Hence it retains, rather than discards, all common likelihood components. The conclusion is
\[
\inf_{\widehat G_n}\sup_{P\in\mathcal M_{\mathrm{tr}}(\gamma;\kappa)}
 \mathbb E_P\!\left[d_H(\widehat G_n,G_P;P)\right]
 \ge c_1(r_n^*)^2
\]
in both branches. By \(\mathrm{def:minimax-risk}\), the left side is \(\mathfrak R_n(1,\gamma;\kappa)\). All constants used in the concentration and packing calculations depend only on the fixed calibrated parameters, the fixed smoothness indices, and the fixed basis orders, not on \(n\). Decreasing \(c_1\) or increasing \(C_1\), if necessary, gives \(0<c_1<C_1<\infty\), completing the proof.